\setlength{\absparsep}{18pt}
\begin{resumo}
    A diretriz para a elaboração é a NBR 6028/21 segundo a ABNT. Resumo no formato informativo contendo: introdução e apresentação da temática. Apresentação dos objetivos geral e específicos. Metodologia. Resultados e discussão e conclusão. Deve ser no impessoal, na terceira pessoa do singular, na voz ativa com no mínimo 150 e máximo de 500 palavras. O resumo não apresenta nenhum tipo de recuo e nenhum parágrafo. Não deve conter citações. Utilizar tamanho da fonte 12 e espaçamento entre linhas 1,5. Logo após o resumo apresentar sequência de palavras que representam as categorias principais do trabalho que são conhecidas como palavras-chave, usar no mínimo três e no máximo cinco palavras, separadas entre si por ponto e vírgula e finalizadas por ponto. Devem ser grafadas com as iniciais em letra minúscula, exceto os substantivos próprios e nomes científicos, que aparecem com inicial maiúscula.
    
    Palavras-chave: resumo; ABNT; pesquisa.
\end{resumo}
